\scanpage{45}
se, if the index theorem is proved for $X$, then $\unsure{A|Q}_{\unsure{\ell|\chi}}(X,\inserted{\unsure{\xi}})=B_{\unsure{\ell|\chi}}(X)$ is again independant of the polarization, and analoguous remark for $C_{\unsure{\ell|\chi}}(X,\inserted{\unsure{\xi}})$.

When speaking about condition $C_{\unsure{\ell|\chi}}(X,\inserted{\unsure{\xi}})$, emphasize at once its stability properties by products (the proof I suggested works indeed) specialization (with possible change of characteristics), hyperplane or more generally linear sections. Give an extra proposition for the relations with the property $A$, via a formal proposition as follows:

\uline{Proposition}\quad Condititions équivalentes sur $X$(verieté \uline{polarisée}):

(i) $C_{\unsure{\ell|\chi}}(X)$

(ii) $C_{\unsure{\ell|\chi}}(Y)$ et $A_{\unsure{\ell|\chi}}(X\mathrm{x}X)^{\deleted{\unsure{0}}}$, \deleted{\unsure{où l'exposant} $^0$}

(ii bis) $C_{\unsure{\ell|\chi}}(Y)$ et $A_{\unsure{\ell|\chi}}(X\mathrm{x}X)^0$, où l'exposant $^0$ signifie qu'on se borne à exprimer la condition $A$ pour l'homomorphisme en dimension critique
\[
H^{2n-2}\longrightarrow H^{2n+2}.
\]

(iii) $C_{\unsure{\ell|\chi}}(Y)$ et $A_{\unsure{\ell|\chi}}(X\mathrm{x}Y)$.

(iii bis) $C_{\unsure{\ell|\chi}}(Y)$ et $A_{\unsure{\ell|\chi}}(X\mathrm{x}Y)^0$, où l'exposant $^0$ signifie qu'on se borne à exprimer la condition $A$ en dimension critique $H^{(2n-1)-1}\longrightarrow H^{(2n-1)+1}$.

(iv) \deleted{Pour tout $\unsure{n|O|0}$} $C_{\unsure{\ell|\chi}}(Y)$, et pour tout $i\leq n-1$, l'homomorphisme naturel $H^i(Y)\mathbin{-}H^i(X)$ inverse à gauche de $\varphi^i:H^i(X)\mathbin{-}H^i(Y)$ (induit par $\Lambda_X\varphi_{\unsure{*}}$) est induit par une classe de correspondance algébrique (induisant ce qu'elle veut sur les autres $H^j(Y)$).

(iv bis) $C_{\unsure{\ell|\chi}}(Y)$, et pour $j\geq n+1$, l'homomorphisme naturel $H^j(X)\longrightarrow H^{j\inserted{\unsure{-2|2}}}(Y$ inverse à droite de $\varphi_{j-2}:H^{j-2}(Y)\mathbin{-}H^j(X)$ \deleted{est} (induit par $\varphi^*\Lambda_X$) est induit par une classe de correspondance algébrique (induisant ce qu'elle veut sur les autres $H^i(X)$).

\uline{Corollaire}\quad \unsure{Cés|Ces} conditions équivalent aussi à
\[
\text{(v)}\qquad A_{\unsure{\ell|\chi}}(X\mathrm{x}X)^0+A_{\unsure{\ell|l}}(Y\mathrm{x}Y)^0+A_{\unsure{\ell|l}}(Z\mathrm{x}Z)^0+\qquad,
\]
où $X\supset Y\supset Z$ est une suite décroissante de sections hyperplanes\inserted{ non sing.}
\mspage{44}
